% SGA 3, Expose VIII, English reconstruction.
% Scope: all of Section 6, stopping immediately before the Section 7 heading.
% Authority: Polo--Gille Exp8-8nov09.pdf, standalone pages 17--22
% (combined-reader pages 633--638).
% The Jacob Reinhold Markdown lineage was used only as a phrasing comparison.
% The Polo--Gille PDF controls prose, formulas, numbering, references, and notes.

% Authority standalone p. 17; original-edition pagination marker 20.
\SGARefTarget{sga3:viii:section:6}\section*{6. 本质自由态射与某些形如 \(\prod_{Y/S}Z/Y\) 的函子的可表性\textsuperscript{(*)}}
\addcontentsline{toc}{section}{6. 本质自由态射与某些函子的可表性}

\begingroup
\renewcommand{\thefootnote}{(*)}
\footnotetext[0]{\SGARefLink{sga3:viii:section:6}{本节}独立于可对角化群理论；其自然位置应在第 VI\textsubscript{B} 讲中。}
\endgroup

\noindent\SGARefTarget{sga3:viii:definition:6:1}\textbf{定义 6.1.}
设 \(f:X\to S\) 为预概形态射。若存在 \(S\) 的仿射开覆盖 \((S_i)\)，对每个 \(i\) 存在在 \(S_i\) 上仿射且忠实平坦的 \(S_i\)-预概形 \(S_i'\)，并存在 \(X_i'=X\times_SS_i'\) 的仿射开子集 \(X'_{ij}\) 构成的覆盖 \((X'_{ij})_j\)，使得对每个 \((i,j)\)，\(X'_{ij}\) 的环为 \(S_i'\) 的环上的自由模，就称 \(f\) 为\emph{本质自由的}，或称 \(X\) \emph{在 \(S\) 上本质自由}。\footnote{编者注：“自由”可换为“投射”；参见下文 \SGARefLink{sga3:viii:remark:6:8}{6.8}。还应将此概念与 \SGARefLink{sga3:viii:bib:rg71}{[RG71]} 中引入并发展的平坦且纯的 \(S\)-预概形概念比较；特别见\emph{同上}，第一部分，3.3.12，以及第二部分，3.1.4.1。}

% Authority standalone p. 18.
\medskip
\noindent\SGARefTarget{sga3:viii:proposition:6:2}\textbf{命题 6.2.}
\begin{enumerate}
  \item[(a)]\SGARefTarget{sga3:viii:proposition:6:2:item:a:01}
  \textit{若 \(X\) 在 \(S\) 上本质自由，则它在 \(S\) 上平坦。当 \(S\) 为 Artin 预概形时，逆命题成立。}

  \item[(b)]\SGARefTarget{sga3:viii:proposition:6:2:item:b:01}
  \textit{若 \(S\) 为域的谱，则每个 \(S\)-预概形都在 \(S\) 上本质自由。}

  \item[(c)]\SGARefTarget{sga3:viii:proposition:6:2:item:c:01}
  \textit{若 \(X\) 在 \(S\) 上本质自由，且 \(S'\to S\) 为基变换态射，则 \(X'=X\times_SS'\) 在 \(S'\) 上本质自由。若 \(S'\to S\) 忠实平坦且拟紧，逆命题也成立。}
\end{enumerate}

证明直接；在证明 \SGARefLink{sga3:viii:proposition:6:2:item:a:01}{(a)} 的逆命题时，利用 Artin 局部环上的平坦模为自由模这一事实。\footnote{编者注：事实上，设 \((A,\mathfrak m)\) 为 Artin 局部环，\(k\) 为其剩余域，\(M\) 为任意 \(A\)-模，\((x_i)_{i\in I}\) 为 \(M\) 的元素，其像构成 \(M/\mathfrak mM\) 在 \(k\) 上的一组基。令 \(F\) 为以 \((e_i)_{i\in I}\) 为基的自由 \(A\)-模，\(\phi:F\to M\) 为由 \(\phi(e_i)=x_i\) 定义的 \(A\)-态射。则 \(Q=\operatorname{Coker}\phi\) 满足 \(Q=\mathfrak mQ\)，从而由于 \(\mathfrak m\) 幂零，有 \(Q=0\)。再假设 \(M\) 在 \(A\) 上平坦。则 \(K=\operatorname{Ker}\phi\) 满足 \(K\otimes_Ak=0\)，即 \(K=\mathfrak mK\)，从而 \(K=0\)。}

\medskip
\noindent\SGARefTarget{sga3:viii:proposition:6:3}\textbf{命题 6.3.}
\textit{设 \(H\) 为可对角化 \(S\)-群预概形（更一般地，在对 \(S\) 的每个仿射开子集作适当的忠实平坦拟紧扩张后变为可对角化；即 \(H\) 为“乘法型”，参见 IX，\SGARefLink{sga3:ix:definition:1:1}{1.1}）。则 \(H\) 在 \(S\) 上本质自由。}

事实上，若 \(H\) 可对角化，则它在 \(S\) 上仿射，且由一个作为 \(\mathcal O_S\)-模自由的代数定义。

引入\SGARefLink{sga3:viii:definition:6:1}{定义 6.1} 的理由在于\SGARefLink{sga3:viii:theorem:6:4}{下述定理}。

% Original-edition pagination marker 21.
\medskip
\noindent\SGARefTarget{sga3:viii:theorem:6:4}\textbf{定理 6.4.}
\textit{设 \(S\) 为预概形，\(Z\) 为在 \(S\) 上本质自由的 \(S\)-预概形，\(Y\) 为 \(Z\) 的闭子预概形。考虑函子}\footnote{编者注：参见 II.1，其中\SGARefLink{sga3:viii:theorem:6:4:display:functor:01}{此函子}记为 \(\prod_{Z/S}Y\)。}
\SGARefTarget{sga3:viii:theorem:6:4:display:functor:01}
\[
  \begin{aligned}
  F=\prod_{Z/S}Y/Z:\;(\mathbf{Sch}/S)^\circ
    &\longrightarrow\mathbf{Set},\\
  F(S')=\Gamma(Y_{S'}/Z_{S'})
    &=
    \begin{cases}
      \varnothing,& Z_{S'}\ne Y_{S'},\\
      \{\operatorname{id}_{Z_{S'}}\},& Z_{S'}=Y_{S'}.
    \end{cases}
  \end{aligned}
\]
\textit{\SGARefLink{sga3:viii:theorem:6:4:display:functor:01}{此函子}可由 \(S\) 的一个闭子预概形表示。}\footnote{编者注：将 \(Z\) 改为 \(S\)。}

\smallskip
\noindent\footnote{编者注：以下扩充了原文；另见 \SGARefLink{sga3:viii:corollary:1:7:3}{1.7.3}。}
首先注意，\(F\) 为关于 \((\mathrm{fpqc})\) 拓扑的层。由于对每个 \(S'\)，\(F(S')\) 或为 \(\varnothing\)，或为 \(\{\mathrm{pt}\}\)，只须验证如下断言：若 \((S_i)\) 为 \(S\) 的开覆盖（分别地，若 \(S'\to S\) 忠实平坦且拟紧），且每个 \(Y_{S_i}\to Z_{S_i}\) 为同构（分别地，\(Y_{S'}\to Z_{S'}\) 为同构），则 \(Y\to Z\) 为同构。在第一种情形中这是显然的，在第二种情形中则由 SGA 1，VIII，5.4 或 EGA IV\textsubscript{2}，2.7.1 得出。

此外，由 SGA 1，VIII，2.1 和 5.5，忠实平坦拟紧态射是由闭浸入箭头组成的纤维范畴的有效下降态射。因此，采用 \SGARefLink{sga3:viii:definition:6:1}{6.1} 的记号，可以归结到 \(S=S_i'\) 的情形。

令 \((Z_j)\) 为 \(Z\) 的仿射开覆盖，使 \(\mathcal O(Z_j)\) 为 \(A=\mathcal O(S)\) 上的自由模，置 \(Y_j=Y\cap Z_j\)，并令
\[
  F_j:(\mathbf{Sch}/S)^\circ\longrightarrow\mathbf{Set}
\]
为由 \((Z_j,Y_j)\) 按照从 \((Z,Y)\) 定义 \(F\) 的同样方式定义的函子。它是终函子的子函子，且显然
\[
  F=\bigcap_jF_j.
\]

% Authority standalone p. 19.
因此只须证明，每个 \(F_j\) 都可由 \(S\) 的一个闭子概形 \(T_j\) 表示；此时 \(F\) 就由闭子概形 \(T=\bigcap_jT_j\) 表示。从而还可假设 \(Z\) 仿射，设为 \(Z=\operatorname{Spec}(B)\)，其中 \(B\) 为自由 \(A\)-模。令 \(J\) 为 \(B\) 中定义 \(Z\) 的闭子预概形 \(Y\) 的子集，\(K\) 为 \(A\) 中由各 \(u_i(J)\subseteq A\) 生成的理想，其中各 \(u_i:B\to A\) 为所选基的坐标线性形式。立即可见，
\[
  T=\mathcal V(K)=\operatorname{Spec}(A/K)
\]
具有所需性质。证明完毕。

\medskip
\noindent\SGARefTarget{sga3:viii:examples:6:5}\textbf{例 6.5.}
给出若干重要例子，它们可归结为 \SGARefLink{sga3:viii:theorem:6:4}{6.4} 所考虑的类型的函子 \(\prod_{Z/S}Y/Z\)，下文将用到它们的可表性判据。这里 \(S\) 表示预概形，\(X,Y,Z,\ldots\) 表示 \(S\) 上的预概形。

\smallskip
\noindent\SGARefTarget{sga3:viii:examples:6:5:item:a:01}\textit{(a)}
假设给定 \(S\)-态射
\[
  q:X\longrightarrow\underline{\operatorname{Hom}}_S(Y,Z),
  \tag{x}\label{sga3:viii:examples:6:5:eq:x:01}
\]
（“\(X\) 作用于 \(Y\)，取值于 \(Z\)”），即态射
\[
  r:X\times_SY\longrightarrow Z.
  \tag{xx}\label{sga3:viii:examples:6:5:eq:xx:01}
\]
设 \(Z'\) 为 \(Z\) 的子预概形。它给出单态射
\[
  \underline{\operatorname{Hom}}_S(Y,Z')
  \longrightarrow
  \underline{\operatorname{Hom}}_S(Y,Z),
\]
使第一个函子成为第二个的子函子。令 \(X'\) 为此子函子在 \SGARefLink{sga3:viii:examples:6:5:eq:x:01}{(x)} 下的逆像。于是 \(X'\) 为 \(X\) 的子函子，其中 \(X'(T)\) 为使
\[
  q(x):Y_T\longrightarrow Z_T
\]
经 \(Z_T'\) 分解的 \(x\in X(T)\) 所成的集合。此函子可作如下描述。置 \(P=X\times_SY\)，令 \(P'\) 为 \(Z'\) 在 \(r:P\to Z\) 下的逆像。则有显然的同构
\[
  X'\simeq\prod_{P/X}P'/P.
  \tag{xxx}\label{sga3:viii:examples:6:5:eq:xxx:01}
\]
由此可知，\textit{若 \(Y\) 在 \(S\) 上本质自由，且 \(Z'\) 在 \(Z\) 中闭，则 \(X\) 的子函子 \(X'\) 可由 \(X\) 的闭子预概形表示。}

\smallskip
\noindent\SGARefTarget{sga3:viii:examples:6:5:item:b:01}\textit{(b)}
假设给定 \(X\) 在 \(Y\) 上取值于 \(Z\) 的两个作用，即两个态射
\[
  q_1,q_2:X\rightrightarrows
  \underline{\operatorname{Hom}}_S(Y,Z),
\]
并置 \(X'=\operatorname{Ker}(q_1,q_2)\)。这是 \(X\) 的子函子，其中 \(X'(T)\) 为使两个态射
\[
  q_1(x),q_2(x):Y_T\rightrightarrows Z_T
\]
相等的 \(x\in X(T)\) 所成的集合。给定 \(q_1,q_2\) 等价于给定态射
\[
  q:X\longrightarrow
  \underline{\operatorname{Hom}}_S(Y,Z\times_SZ),
\]
也等价于给定态射
\[
  r:X\times_SY\longrightarrow Z\times_SZ.
\]
置 \(U=Z\times_SZ\)，令 \(U'\) 为 \(Z\times_SZ\) 的对角子预概形。则 \(X'\) 为子函子
\[
  \underline{\operatorname{Hom}}_S(Y,U')
  \longrightarrow
  \underline{\operatorname{Hom}}_S(Y,U).
\]
在 \(q\) 下的逆像。因此它可写成 \SGARefLink{sga3:viii:examples:6:5:eq:xxx:01}{(xxx)} 的形式，其中 \(P=X\times_SY\)，\(P'\) 为对角在 \(r\) 下的逆像，即
\[
  X\times_SY
  \mathop{\rightrightarrows}^{r_1}_{r_2}
  Z.
\]
的核。因此处于 \SGARefLink{sga3:viii:examples:6:5:item:a:01}{(a)} 的情形。

% Authority standalone p. 20; original-edition pagination marker 24.
从而，\textit{若 \(Y\) 在 \(S\) 上本质自由，且 \(Z\) 在 \(S\) 上分离，则 \(X\) 的子函子 \(X'\) 可由 \(X\) 的闭子预概形表示。}

\smallskip
\noindent\SGARefTarget{sga3:viii:examples:6:5:item:c:01}\textit{(c)}
假设给定态射
\SGARefTarget{sga3:viii:examples:6:5:item:c:display:q:01}\[
  q:X\longrightarrow\underline{\operatorname{Hom}}_S(Y,Y);
\]
即“\(X\) 作用于 \(Y\)”。令 \(X'\) 为\SGARefLink{sga3:viii:examples:6:5:item:c:display:q:01}{此态射}的“核”，即 \(X\) 中满足如下条件的子函子：\(X'(T)\) 为使 \(q(x):Y_T\to Y_T\) 为恒等映射的 \(x\in X(T)\) 所成的集合。引入第二个态射
\[
  q':X\longrightarrow\underline{\operatorname{Hom}}_S(Y,Y)
\]
使“\(X\) 平凡作用于 \(Y\)”，即可看出此函子属于 \SGARefLink{sga3:viii:examples:6:5:item:b:01}{(b)} 的情形。因此，\textit{若 \(Y\) 在 \(S\) 上本质自由且分离，则 \(q\) 的核子函子可由 \(X\) 的闭子预概形表示。}

\smallskip
\noindent\SGARefTarget{sga3:viii:examples:6:5:item:d:01}\textit{(d)}
在 \SGARefLink{sga3:viii:examples:6:5:item:c:01}{(c)} 的假设下，考虑 \(Y\) 中“在 \(X\) 下不变的”子函子 \(Y'\)。于是 \(Y'(T)\) 为满足下述条件的 \(y\in Y(T)\) 所成的集合：相应态射 \(\bar q(y):X_T\to Y_T\) 为“取值为 \(y\) 的 \(T\)-常值态射”。如 \SGARefLink{sga3:viii:examples:6:5:item:c:01}{(c)} 引入 \(q'\)，并用 \(\bar q\) 和 \(\bar q'\) 表示相应态射，
\[
  \bar q,\bar q':Y\rightrightarrows
  \underline{\operatorname{Hom}}_S(X,Y),
\]
便可见 \(Y'=\operatorname{Ker}(\bar q,\bar q')\)。因此它仍属于 \SGARefLink{sga3:viii:examples:6:5:item:b:01}{(b)} 的情形，但将 \(X\) 和 \(Y\) 的角色互换，并取 \(Z=Y\)。从而，\textit{若 \(X\) 在 \(S\) 上本质自由，且 \(Y\) 在 \(S\) 上分离，则 \(Y\) 中在 \(X\) 下不变的子函子 \(Y'\) 可由 \(Y\) 的闭子预概形表示。}

\smallskip
\noindent\SGARefTarget{sga3:viii:examples:6:5:item:e:01}\textit{(e)}
\SGARefLink{sga3:viii:examples:6:5}{上述各例}所描述的构造在群论中特别常见。因此，当 \(S\) 上的群预概形 \(G\) 作用于 \(S\)-预概形 \(X\) 时，
\[
  q:G\longrightarrow\underline{\operatorname{Aut}}_S(X),
\]
由 \SGARefLink{sga3:viii:examples:6:5:item:c:01}{(c)}，只要 \(X\) 在 \(S\) 上本质自由且分离，\(q\) 的核，即“\(G\) 中作用平凡的子群”，就是 \(G\) 的闭子预概形。由 \SGARefLink{sga3:viii:examples:6:5:item:d:01}{(d)}，只要 \(G\) 在 \(S\) 上本质自由且 \(X\) 在 \(S\) 上分离，不变子对象 \(X^G\) 就是 \(X\) 的闭子预概形。\footnote{编者注：将 \(S_X\) 改为 \(S\)。}

设 \(Y,Z\) 为 \(X\) 的子预概形，考虑 \(G\) 的子函子 \(\underline{\operatorname{Transp}}_G(Y,Z)\)，即“将 \(Y\) 输送到 \(Z\) 中的输送子”。它在 \(S\)-预概形 \(T\) 中取值的点，是使相应的 \(X_T\) 的自同构满足 \(g(Y_T)\subseteq Z_T\) 的 \(g\in G(T)\)，即诱导的态射 \(Y_T\to X_T\) 经 \(Z_T\) 分解。因此，\textit{若 \(Y\) 在 \(S\) 上本质自由，且 \(Z\) 在 \(X\) 中闭，则由 \SGARefLink{sga3:viii:examples:6:5:item:a:01}{(a)}，\(\underline{\operatorname{Transp}}_G(Y,Z)\) 是 \(G\) 的闭子预概形。}

还可考虑将 \(Y\) 输送到 \(Z\) 的\emph{严格输送子}，\footnote{编者注：它记为 \(\underline{\operatorname{Transp}}^{\mathrm{str}}_G(Y,Z)\)。}其 \(T\)-值点为使 \(g(Y_T)=Z_T\) 的 \(g\in G(T)\)。它恰为
\[
  \underline{\operatorname{Transp}}_G(Y,Z)
  \cap
  \sigma\bigl(\underline{\operatorname{Transp}}_G(Z,Y)\bigr),
\]
其中 \(\sigma\) 为 \(G\) 的对称。因此，若 \(Y\) 和 \(Z\) 在 \(S\) 上本质自由且在 \(X\) 中闭，则将 \(Y\) 输送到 \(Z\) 的严格输送子是 \(G\) 的闭子预概形。

% Authority standalone p. 21; original-edition pagination marker 25.
一个重要情形是 \(X=G\)，其中 \(G\) 通过内自同构作用于自身。若 \(H\) 为 \(G\) 的子预概形，则将 \(H\) 输送到自身的严格输送子也称为 \emph{\(H\) 在 \(G\) 中的正规化子}，记作 \(\underline{\operatorname{Norm}}_G H\)。因此，\textit{若 \(H\) 为 \(G\) 的闭子群预概形，且在 \(S\) 上本质自由，则 \(\underline{\operatorname{Norm}}_G H\) 可由 \(G\) 的闭子群预概形表示。}

最后，设 \(Z\) 为 \(G\) 的子预概形。它在 \(G\) 中的\emph{中心化子} \(\underline{\operatorname{Centr}}_G(Z)\) 为用 \SGARefLink{sga3:viii:examples:6:5:item:d:01}{(d)} 的方法定义的 \(G\) 的群子函子，这里通过 \(G\) 的作用所诱导的作用来理解“\(Z\) 作用于 \(G\)”。于是，\textit{若 \(Z\) 在 \(S\) 上本质自由，且 \(G\) 在 \(S\) 上分离，则 \(\underline{\operatorname{Centr}}_G(Z)\) 为 \(G\) 的闭子群预概形。}特别地，\textit{若 \(G\) 在 \(S\) 上本质自由且分离，则 \(G\) 的中心 \(C\)，也就是 \(\underline{\operatorname{Centr}}_G(G)\)，为 \(G\) 的闭子群预概形。}

当 \(S\) 为域的谱时，\SGARefLink{sga3:viii:proposition:6:2:item:b:01}{6.2(b)}\emph{〔原文印作 6.3(b)〕} 表明，上述各例 \SGARefLink{sga3:viii:examples:6:5:item:a:01}{(a)}--\SGARefLink{sga3:viii:examples:6:5:item:e:01}{(e)} 中“本质自由”的假设自动满足；只剩下分离性假设。回顾域上的群预概形必分离，例如可得：

\medskip
\noindent\SGARefTarget{sga3:viii:corollary:6:7}\textbf{推论 6.7.}
\footnote{编者注：保留了原文编号；没有 6.6。}
\textit{设 \(G\) 为域 \(k\) 上的群预概形。则：}
\begin{list}{--}{\setlength{\leftmargin}{1.8em}\setlength{\itemsep}{0.35em}}
  \item\SGARefTarget{sga3:viii:corollary:6:7:item:1:01}
  \textit{对 \(G\) 的每个子预概形 \(Z\)，\(Z\) 在 \(G\) 中的中心化子为 \(G\) 的闭子群预概形。特别地，这对 \(G\) 的中心 \(\underline{\operatorname{Centr}}_G(G)\) 成立。}

  \item\SGARefTarget{sga3:viii:corollary:6:7:item:2:01}
  \textit{更一般地，若 \(u,v:X\to G\) 为预概形态射，则 \(\underline{\operatorname{Transp}}_G(u,v)\) 可由 \(G\) 的闭子预概形表示。}

  \item\SGARefTarget{sga3:viii:corollary:6:7:item:3:01}
  \textit{对 \(G\) 的两个子预概形 \(Y,Z\)，若 \(Z\) 为闭的，则 \(\underline{\operatorname{Transp}}_G(Y,Z)\) 为 \(G\) 的闭子预概形。若 \(Y\) 也为闭的，则\SGARefLink{sga3:viii:corollary:6:7:item:3:01}{同一结论}对 \(\underline{\operatorname{Transp}}^{\mathrm{str}}_G(Y,Z)\) 成立。}

  \item\SGARefTarget{sga3:viii:corollary:6:7:item:4:01}
  \textit{对 \(G\) 的每个子群预概形}\footnote{编者注：事实上，在域 \(k\) 上，\(G\) 的每个子群预概形都是闭的；参见第 VI\textsubscript{A} 讲，0.6.1。}
  \textit{\(H\)，正规化子 \(\underline{\operatorname{Norm}}_G(H)\) 为 \(G\) 的闭子群预概形。}
\end{list}

\medskip
\noindent\SGARefTarget{sga3:viii:remark:6:8}\textbf{注 6.8.}
\footnote{编者注：以下扩充了原文；特别增加了\SGARefLink{sga3:viii:lemma:6:8:1}{引理 6.8.1}。}
设 \(A\) 为交换环，\(M\) 为 \(A\)-模，并置
\[
  M^\vee=\operatorname{Hom}_A(M,A).
\]
在 \(A\)-模 \(\operatorname{End}_A(M)\) 上赋予离散逐点收敛拓扑。于是，零点的一个邻域基由如下 \(A\)-子模组成，其中 \(n\in\mathbb N\) 且 \(m_1,\ldots,m_n\in M\)：
\SGARefTarget{sga3:viii:remark:6:8:display:neighborhood:01}
\[
  K(m_1,\ldots,m_n)
  =
  \{u\in\operatorname{End}_A(M)\mid
  u(m_i)=0\ \text{对 \(i=1,\ldots,n\)}\}.
\]
若典范态射
\[
  \Theta:M\otimes_AM^\vee\longrightarrow\operatorname{End}_A(M)
\]
的像的闭包包含 \(\operatorname{id}_M\)，即\SGARefLink{sga3:viii:remark:6:8:eq:star:01}{下列条件}成立，就称 \(M\) 为\emph{拟自由 \(A\)-模}：

% Authority standalone p. 22.
\[
  \left\{
  \begin{aligned}
    &\text{对任意 \(m_1,\ldots,m_n\in M\)，存在
      \(x_1,\ldots,x_r\in M\) 和 \(f_1,\ldots,f_r\in M^\vee\)}\\
    &\text{使得 }\quad
    m_i
    =
    \Theta\!\left(\sum_{s=1}^r x_s\otimes f_s\right)(m_i)
    =
    \sum_{s=1}^r f_s(m_i)x_s
    \quad\text{对 \(i=1,\ldots,n\)。}
  \end{aligned}
  \right.
  \tag{*}\label{sga3:viii:remark:6:8:eq:star:01}
\]
\noindent（此时 \(\operatorname{Im}\Theta\) 在 \(\operatorname{End}_A(M)\) 中稠密，因为对每个 \(u\in\operatorname{End}_A(M)\)，有
\(
  u(m_i)
  =
  \sum_{s=1}^r f_s(m_i)u(x_s)
  =
  \Theta\!\left(\sum_{s=1}^r u(x_s)\otimes f_s\right)(m_i).
\)）

首先注意，\SGARefLink{sga3:viii:remark:6:8:eq:star:01}{此性质}在基变换下稳定。事实上，设 \(\phi:A\to A'\) 为环同态，置 \(M'=M\otimes_AA'\)，并取 \(m_1',\ldots,m_n'\in M'\)。写成
\[
  m_i'=\sum_jm_{ij}\otimes b_{ij},
  \qquad m_{ij}\in M,\quad b_{ij}\in A'.
\]
由假设，存在 \(x_1,\ldots,x_r\in M\) 和 \(f_1,\ldots,f_r\in\operatorname{Hom}_A(M,A)\)，使
\[
  m_{ij}=\sum_sx_sf_s(m_{ij})
  \qquad\text{对所有 \(i,j\)。}
\]
用 \(\phi\circ f_s\) 表示 \(f_s\) 在 \[
  \operatorname{Hom}_A(M,A')
  =
  \operatorname{Hom}_{A'}(M',A').
\] 中的像。
则对 \(i=1,\ldots,n\)，
\SGARefTarget{sga3:viii:remark:6:8:display:base-change:01}
\[
  \begin{aligned}
  \left(\sum_sx_s\otimes\phi\circ f_s\right)(m_i')
    &=
    \sum_{s,j}x_s\otimes\phi(f_s(m_{ij}))\,b_{ij}\\
    &=
    \sum_j\left(\sum_sx_sf_s(m_{ij})\right)\otimes b_{ij}\\
    &=m_i'.
  \end{aligned}
\]
这证明了 \(M'\) 在 \(A'\) 上拟自由。

此外，每个投射 \(A\)-模 \(P\) 都是拟自由的。事实上，存在 \(A\)-态射
\[
  A^{(I)}\xrightarrow{\ \pi\ }P
  \xrightarrow{\ \tau\ }A^{(I)}
  \qquad\text{使得}\qquad
  \pi\circ\tau=\operatorname{id}_P.
\]
令 \((e_i)\) 为 \(A^{(I)}\) 的典范基，并置 \(f_i=e_i^*\circ\tau\)，它是 \(P\) 上的线性形式。给定 \(m_1,\ldots,m_n\in P\)，存在 \(I\) 的有限子集 \(J\)，使得
\[
  m_k=\sum_{i\in J}f_i(m_k)\pi(e_i)
  \qquad\text{对 \(k=1,\ldots,n\)。}
\]

若在\SGARefLink{sga3:viii:definition:6:1}{定义 6.1} 的陈述中将“自由”换为“投射”，或更一般地换为“拟自由”，定理 \SGARefLink{sga3:viii:theorem:6:4}{6.4} 仍成立。按照 \SGARefLink{sga3:viii:theorem:6:4}{6.4} 的证明进行，即归结到\SGARefLink{sga3:viii:lemma:6:8:1}{下述引理}。

\medskip
\noindent\SGARefTarget{sga3:viii:lemma:6:8:1}\textbf{引理 6.8.1.}
\textit{设 \(M\) 为拟自由 \(A\)-模，\(N\) 为子模，\(F\) 为协变函子}
\[
  F:(\text{\(A\)-代数})\longrightarrow\mathbf{Set}
\]
\textit{使得当 \(M_B=(M/N)_B\) 时 \(F(B)=\{\mathrm{pt}\}\)，否则 \(F(B)=\varnothing\)。则存在 \(A\) 的理想 \(K\)，使得 \(F(B)=\{\mathrm{pt}\}\) 当且仅当态射 \(A\to B\) 经 \(A/K\) 分解；即有关于 \(B\) 函子性的同构}
\SGARefTarget{sga3:viii:lemma:6:8:1:display:representing:01}
\[
  F(B)\simeq\operatorname{Hom}_{A\text{-alg.}}(A/K,B).
\]

\smallskip
\noindent\textit{证明。}
令 \((n_j)\) 为 \(N\) 的生成系，\(F_j\) 为对应于子模 \(An_j\) 的终函子 \(e\) 的子函子，其中对每个 \(B\) 有 \(e(B)=\{\mathrm{pt}\}\)。\(F(B)=\{\mathrm{pt}\}\) 当且仅当每个 \(n_j\) 在 \(M_B\) 中的像为零。因此 \(F\) 为各函子 \(F_j\) 的交，从而归结到 \(N\) 由单个元素 \(n\) 生成的情形。

设 \(B\) 为 \(A\)-代数，结构映射为 \(\phi:A\to B\)。若 \(n\) 在 \(M_B\) 中的像 \(n\otimes1\) 为零，则对每个 \(f\in M^\vee=\operatorname{Hom}_A(M,A)\)，有
\[
  0=f(n)\otimes1=\phi(f(n)).
\]
反过来，由于 \(M\) 拟自由，存在 \(x_1,\ldots,x_r\in M\) 和 \(f_1,\ldots,f_r\in M^\vee\)，使
\[
  n=\sum_sf_s(n)x_s,
\]
从而
\[
  n\otimes1=\sum_sx_s\otimes\phi(f_s(n)).
\]
由此可知，\(n\otimes1=0\) 当且仅当 \(\phi\) 经 \(A/K\) 分解，其中 \(K\) 为由各 \(f_s(n)\)（\(s=1,\ldots,r\)）生成的理想。引理得证。
